\section{Coordinate factors and balanced compatibility separators}
\label{sec:euler-coordinate-factors}

This is a prospective algebraic search rule.  It does not identify geometric
surface components, does not assert that the maintained recognizer already
uses the rule, and does not supply a general quasi-polynomial bound.
Fix an allowed Q mask $J$, write $n=|J|$, and let $C$ be the restriction of
the exact source cycle matrix to those columns.  All evaluations below use
the same full-source canonical Euler function
\[
 K_J=\{q\ge0:Cq=0\},\qquad
 F(q)=\max_{\alpha\in\mathcal A}c_\alpha^{\mathsf T}q,
 \qquad A=|\mathcal A|,
\]
with omitted coordinates extended by zero.  In particular $F$ is positively
homogeneous and subadditive.  Coordinate restriction preserves quadrilateral
compatibility.

\begin{proposition}[An exact coordinate-product certificate]
Let $B_1,\ldots,B_s$ be a partition of $J$ into nonempty sets and put
$C_i=C[:,B_i]$.  The following conditions are equivalent:
\[
 \ker C=\bigoplus_i\ker C_i,
 \qquad
 \operatorname{rank}C=\sum_i\operatorname{rank}C_i,
\]
where each summand is extended by zero.  Under these conditions,
$K_J=\prod_i K_{B_i}$ in disjoint coordinates.
\end{proposition}
\begin{proof}
The displayed direct sum is always contained in $\ker C$.  Its dimension
is $n-\sum_i\operatorname{rank}C_i$, whereas the dimension of $\ker C$ is
$n-\operatorname{rank}C$.  Dimension equality proves the equivalence.
Intersecting the signed-space equality with the nonnegative orthant proves
the cone product.
\end{proof}

A polynomial exact construction chooses a column basis $I$, expresses every
nonbasis column as $C_j=\sum_{i\in I}a_{ij}C_i$, and joins the indices in the
support of each fundamental relation
$e_j-\sum_i a_{ij}e_i$.  These relations span the kernel.  Each support is a
minimal dependence, so their union components form the finest coordinate
partition: the linear-matroid components.  Zero columns are singleton loops;
a singleton coloop has zero kernel.  An independent consumer need not trust
this construction.  It reconstructs the source matrix and checks the literal
partition and the exact rank equality above.  Rational elimination has
polynomial bit complexity in the source encoding length.

\begin{proposition}[Completeness after implicit-zero deletion]
If $K_J$ contains a vector strictly positive on $J$, its literal coordinate
product partitions are exactly the rank-additive partitions of $C$.
In general, let $J^+$ contain the coordinates positive in at least one
vector of $K_J$.  After deleting $J\setminus J^+$, the matroid components
give the finest coordinate product partition of the remaining cone.
\end{proposition}
\begin{proof}
If $u>0$ lies in $K_J$ and $v\in\ker C$, then $u+\varepsilon v\ge0$ for
sufficiently small positive $\varepsilon$.  Both this vector and $u$ are
in $K_J$, so $\operatorname{span}K_J=\ker C$.  A coordinate product of
the cone gives a direct sum of these spans.  Each factor is its intersection
with that coordinate axis, since the other factors contain zero.
In each block the restricted
positive vector proves that its cone spans the full restricted kernel.
Thus a cone product forces the signed-kernel product and rank additivity;
the converse was proved above.  For $J^+$, summing one nonnegative witness
per positive coordinate gives a vector strictly positive on every retained
coordinate.  All deleted coordinates vanish identically, so the same
argument applies.
\end{proof}

The set $J^+$ can be found by polynomially many normalized coordinate LPs:
test whether the maximum of $q_j$ subject to $Cq=0$, $q\ge0$ and
$\mathbf1^{\mathsf T}q\le1$ is positive, for each $j$.
A nonpositive dual certifies deletion; a positive answer retains the
coordinate.  A capped query authorizes no deletion.  With complete polynomial
LPs this gives a polynomial construction, but it may cost more than the
rank-only test.  No such new preprocessing is installed here.  These are
products of the formal nonnegative matching cone; compatibility and canonical
Euler impose additional conditions.

\begin{proposition}[Positive search and negative composition]
There is a compatible $q\in K_J$ with $F(q)>0$ if and only if some factor
contains a compatible $q_i\in K_{B_i}$ with $F(q_i)>0$.
\end{proposition}
\begin{proof}
A factor witness extends by zero to a source witness.  Conversely, write a
compatible source vector as $q=\sum_i q_i$ using its literal restrictions to
the factors.  Each restriction is compatible and belongs to its factor cone.
Subadditivity gives $0<F(q)\le\sum_iF(q_i)$, so one summand is positive.
This is a formal-coordinate argument; no assertion about connected normal
surface components is needed.
\end{proof}

Consequently complete negative proofs for all factor masks imply negativity
of the parent mask.  A capped factor does not meet this requirement.  The
existing supplied-source search already accepts each such mask, but a new
outer certificate would have to bind the source, parent mask, rank-certified
partition, and all complete child proofs.  Concatenating child packets does
not make an ordinary branch-tree certificate.  Original STANDARD validation
and independent crushing replay still apply to every returned positive point.

\begin{proposition}[The resolving number adds across coordinate factors]
Use the full-source function $F$ to define each factor's incompatible positive
root-ray conflict hypergraph $\mathcal H_i$, with transversal number $h_i$.
Then the parent resolving number satisfies $h_J=\sum_i h_i$.
\end{proposition}
\begin{proof}
An extreme ray of a product has only one nonzero factor, since two nonzero
factor restrictions give a nontrivial nonnegative decomposition.  Conversely,
a factor extreme ray is extreme in the parent: a nonnegative decomposition
must vanish outside its factor.  Its canonical Euler value and complete
conflict set are unchanged.  Thus the parent hypergraph is the union of the
$\mathcal H_i$.

These hypergraphs have disjoint conflict-bearing tetrahedron sets.  A conflict
needs at least two Q columns in its factor.  Two distinct factors could conflict
at the same tetrahedron only if they contained at least four of its three Q
columns.  Therefore hitting sets for the factor hypergraphs may be chosen
independently and their minimum cardinalities add.  Compatible positive rays
are excluded in every definition; no prior negativity assumption is used.
\end{proof}

For the full source, the three Q columns of each tetrahedron sum to zero,
by the local block-ones identity of Section~\ref{sec:conditional-type-matching}.
To see this in the original matching system, put one in its three Q
coordinates and minus one in its four triangle coordinates, with all
other coordinates zero.  Each face arc is then zero, so this signed
STANDARD vector satisfies every matching equation.  Its Q projection
is the asserted three-column relation.  Restricting this relation to each
matroid component still gives a relation.  Any component containing a nonzero
column from the triple must therefore contain at least two columns of that
triple; there cannot be two such components among three columns.  Hence all
its nonzero columns belong to one matroid component.  Zero columns remain
singleton loop factors with compatible unit rays.  The full all-ones vector
lies in the kernel, so its restriction gives a strictly positive kernel vector in every factor.
This proves nonnegative feasibility, not positive canonical Euler, and does
not classify a factor as negative or positive.

Let $t_i$ count tetrahedra represented by $B_i$, so $\sum_i t_i\le3t$.
Running the existing subset search independently on the factors, with shared
operational accounting, has conditional bit bound
\[
 \sum_i A(3t_i)^{O(h_i)}\operatorname{poly}(t,L).
\]
With $\delta_i$ the maximum of $1$ and all conflict widths encountered in
factor $i$ (so an execution with no conflicts has $\delta_i=1$), the
conflict-directed version of Section~\ref{sec:conflict-candidates} gives
\[
 O\!\left(A\sum_i(3\delta_i)^{h_i}\operatorname{poly}(t,L)\right).
\]
These sums include each factor's unsuccessful candidates and depth restarts.
They can replace exponential dependence on $\sum_i h_i$ by dependence on the
largest factor parameter.  They provide no improvement when the source has
one nontrivial factor.  For authenticated one-vertex stages, $A\le4t$; a full
recognizer would also sum over its polynomially many strict descent stages.

\subsection{A sufficient balanced-separator hypothesis}
Fix a constant base size $m_0$.  Suppose every factor mask with $m>m_0$
columns admits at most $k(t)$ ambiguous tetrahedra such that every singleton
assignment there leaves only coordinate factors with at most $\alpha m$
columns, for a fixed $\alpha<1$.  An assignment retains one currently allowed
Q type at each chosen tetrahedron and deletes the other types there.
These assignments cover every compatible vector, including vectors that
vanish in a chosen tetrahedron.  Call these tetrahedra a balanced compatibility
separator.  There are at most $3^{k(t)}$ assignments, and within each assignment
the child factor sizes sum to at most $m$.  Masks of size at most $m_0$ can be
handled by a constant number of compatible-sector searches, each using at
most $A+1$ LP queries.  A factor path has $O(\log(3t))$ levels.  The size is the number of
retained Q columns, not the number of tetrahedra, and the same global
nondecreasing bound $k(t)$ is required at every descendant.  Thus recursive
branching followed by additive factor composition has at most
$t3^{O(k(t)\log t)}$ factor nodes.  Exhaustive discovery of each separator
costs $(3t)^{O(k(t))}$ polynomial rank work, and fits the overall conditional
bound
\[
 A\,t^{O(k(t))}\operatorname{poly}(t,L).
\]
A polylogarithmic separator theorem would therefore suffice for
quasi-polynomial search without a small whole-source resolving number.
No such theorem is established for the native triangulations.  The condition
is an exact possible next target, and requires a mixed split/factor checker;
small conflicts or short negative explanations alone do not establish it.


To lift this source-stage bound to recognition, the initial triangulation
size must be polynomial in the diagram length; every authenticated successor
must have polynomial-size exact source algebra; the anchor count must remain
polynomially bounded, as it does with $A\le4t$ at the maintained one-vertex
search stages; and the separator hypothesis must hold at every descendant
factor at every stage.  With polynomial preprocessing, independent replay,
primitive basic-ray recovery, and at most
$O(t_0)$ strict-decrease stages from an initial size $t_0$, their sum is still
quasi-polynomial for polylogarithmic $k(t_0)$.  A property of one full root
matrix is insufficient.  Failing the algebraic separator test also does not
rule out immediate LP-negative children, which a complete search can accept
without further decomposition.

\subsection{Exact finite probes and their limitation}

Two separate algebra-only probes use the labelled geometric-arm endpoints
of Section~\ref{sec:geometric-descent-crushing}.  Every original diagram and
complete or partial chain is authenticated again with discovery disabled.
The empty-diagram endpoint is terminal and supplies no Q matrix.  Each
of the other six endpoints has exactly one vertex, lying on the boundary,
and two boundary faces.
The independent verifier's full-source model supplies the matrices; no LP,
recognition search, ray discovery, or timing experiment is run.

The first probe checks an exact row transformation to reduced echelon form,
spanning of the original row space, complete fundamental-circuit kernel
bases, and component ranks in reverse column order.  Every matrix has
one matroid component, no zero columns, no coloops, rank $t-1$, and
nullity $2t+1$.  Component indicator vectors are checked directly against
original matching and canonical Euler.  Such nonnegative witnesses do not
classify positive canonical Euler existence.

The second probe retains every one-tetrahedron singleton assignment:
$3(7+10+16+11+17+31)=276$ masks.  It stores the retained-column map and
complete component data against the original matrix, rather than replacing
that matrix with an unbound projection.  A separate replay reconstructs all
masks, checks complete kernel bases and exact rank additivity, and recomputes
the finest partitions using a different elimination order.

\begin{center}
\begin{tabular}{lrrrr}
\toprule
Source & $t$ & Full Q columns & Singleton masks & Largest remaining factor\\
\midrule
Trefoil & 7 & 21 & 21 & 19\\
Figure eight & 10 & 30 & 30 & 28\\
$TT$ & 16 & 48 & 48 & 46\\
$T(2,5)$ & 11 & 33 & 33 & 31\\
$T(2,7)$ & 17 & 51 & 51 & 49\\
$W(3,4)$ & 31 & 93 & 93 & 91\\
\bottomrule
\end{tabular}
\end{center}

Every one of the 276 masks still has one matroid component, with no zero
columns or coloops.  Its size is exactly $3t-2$, and it includes a retained
column from every original tetrahedron.  All tetrahedra tie under the
best worst-case measure over their three assignments.  There is no algebraic
reduction beyond the two types removed by the assignment itself.
In particular none of these one-tetrahedron choices attains a two-thirds
balance.  These finite results neither rule out larger separators nor prove
an asymptotic lower bound.  Signed-kernel connectivity alone does not certify
nonnegative feasibility or establish that the Euler queries are difficult.
Separately, Section~\ref{sec:conditional-type-matching} supplies original
STANDARD nonnegative matching witnesses for all $276$ singleton masks on
these same labelled source tables.  Those witnesses allow incompatible Q
types in the other tetrahedra and do not establish positive canonical Euler.
They do establish full feasible support in each tested singleton mask.
Indeed its retained type has a positive coordinate in the supplied witness;
adding a block-ones kernel vector in every other tetrahedron makes all
remaining Q coordinates strictly positive while preserving the singleton
restriction.  The full roots already have the all-ones vector.  By the
completeness proposition, the connected matroids therefore rule out every
nontrivial literal coordinate product of these six full and $276$ restricted
nonnegative Q cones.  This conclusion uses the prior matching witnesses and
the proved local identity, not signed-kernel connectivity alone or a new LP
experiment.  It does not rule out other search decompositions.

Root-only factorization and a one-tetrahedron split therefore offer no
factor-search benefit on this fixed corpus.  No factor producer or new
negative certificate schema is installed.  The retained scripts, full
records, logs, and portable independent replay are under
\path{data/coordinate-factor-*}; their source bindings distinguish the
historical geometric discovery from the fresh algebra-only work.
The portable replay authenticates all seven chains, six full matrices and
$276$ restricted matrices without LP or recognition discovery.  Its first
inventory attempt stopped before these checks because a repository-relative
Git path was used from a project subdirectory.  The failed source and outputs
are retained; only repository-root resolution changed before the successful
replay.  No experiment record or native algorithm changed in that repair.

\subsection{A precise limitation of a connected-sum construction}

A different route to understanding a large resolving number would embed
many old positive incompatible rays in one actual knot exterior.  The
following limited transport statement removes one of its obligations.
Suppose an old triangulated piece is retained literally in a larger finite
triangulation by attaching only former boundary faces, so every old internal
matching equation remains.  Suppose an old nonzero Q extreme ray $r$ admits
a global matching extension whose Q vector is $(r,0)$ on old and new
tetrahedra.  Then this global Q ray is extreme.  Indeed a nonnegative
decomposition of $(r,0)$ vanishes on all new Q coordinates.  Restricting
matching lifts of its summands to the old tetrahedra gives old Q-kernel
vectors.  Old extremality makes both proportional to $r$, and their zero
extensions are therefore proportional globally.

This statement concerns strict Q-support locality.  Merely keeping the
\emph{conflicts} in the old piece while adding compatible nonzero Q
coordinates elsewhere does not give the same argument: the extension fiber
could contain new directions.  Nor does extremality alone preserve canonical
Euler positivity.  A connected-sum family still needs a genuine gluing,
matching triangle traces, canonical Euler control, growth relative to source
size, and persistence through the actual preprocessing policy.

There is a direct obstruction to the most literal proposed port.  An embedded
annulus that is a boundary-face subcomplex has two disjoint nonempty boundary
circles.  Each circle contains a triangulation vertex.  A one-vertex torus
boundary cannot contain these two disjoint circles as subcomplexes, because
both would contain its sole vertex.  Thus the six present endpoints do not
already carry such an embedded annular boundary-face port.  The Q-zero
$2\leftrightarrow3$ transport of Section~\ref{sec:q-zero-pachner-transport}
preserves the vertices and cannot by itself create one.  Boundary subdivision
with a separately proved witness map, or cutting along a normal annulus with
controlled trace and Euler transport, would be additional work.  The observed
$h=1$ and $h=2$ cases do not replace those construction obligations.
